\section{Methodology}
\subsection{Judgment Prediction}
The judgment prediction task involves both binary classification (e.g., favoring or opposing a party) and ternary classification (e.g., fully accepted, partially accepted, or rejected). We employ two distinct approaches for this purpose: Language Model-based and LLM-based strategies. These approaches are carefully designed to handle the complexity and diversity of Indian legal documents, ensuring robust performance across various scenarios.

\subsubsection{Language Model based}
In our approach, we utilized several language models, including InLegalBERT, InCaseLaw \cite{paul-2022-pretraining}, and XLNet (large) \cite{yang2019xlnet}, as baselines for binary and ternary classification. Due to the length constraints of complete judgments, which exceed the token capacity of these models, we adopted a chunking strategy. Each document was divided into 512-token chunks using a moving window approach with a 100-token overlap to preserve textual context.

\subsubsection{Large Language Model based}
To utilize LLMs in prediction, we employed two strategies: one involving only prediction instructions and the other prediction with explanation instructions. We used two methods to get predictions from \texttt{\namel} after CPT and CPT followed by SFT. We followed the prompts and instruction-tuning approaches published by \cite{vats-etal-2023-llms} and \cite{nigam2024legaljudgmentreimaginedpredex} in a few-shot setup, and used the PredEx training data for instruction-tuning.

\subsection{Judgment Prediction with Explanation}
For this task, we employed the same LLMs with settings similar to the Judgment Prediction task, but with modified instructions focusing on both prediction and explanation.

\subsection{Prompts Used}
For both task inferences we utilized prompts from \cite{vats-etal-2023-llms}. These prompts, which include a case description and a gold standard prediction label, guide the LLM to generate judicial decisions. The details of these prompts can be found in Table \ref{tab:judgment_prediction_prompts_few} in the Appendix. In addition, for instruction tuning, we adopted prompts from \cite{nigam2024legaljudgmentreimaginedpredex} for prediction tasks, as listed in Table \ref{tab:judgment_prediction_prompts_zero} in the Appendix. The relevant details and examples of the generated predictions and explanations are also available in the referenced tables in the Appendix.

\subsection{Instruction-Set}
For both judgment prediction and explanation, we used 16 instruction sets correspondingly published by \cite{nigam2024legaljudgmentreimaginedpredex}. For a comprehensive view of all instruction sets which was randomly given to the model for tuning, we have included the full list in Table~\ref{Instruction-sets} in the Appendix of this paper.